% called by main.tex
%
\chapter{Reproducibilidad}
\label{ch::reproducibilidad}

\section{Enlace a repositorio GitHub}

El código, los registros de resultados y los artefactos públicos de auditoría se mantienen en
un repositorio:

\begin{center}
\url{https://github.com/marcmaldonadolorca/polymarket-btc-microstructure}
\end{center}

\subsection{Estructura del repositorio}

El repositorio se organiza en los siguientes directorios:

\begin{description}[align=right,labelwidth=2.6cm]
\item [\texttt{config/}] Umbrales congelados de la política y configuraciones del arco.
\item [\texttt{docs/}] Documentos canónicos del proceso, renumerados de forma secuencial.
\item [\texttt{notebooks/}] Cuadernos pedagógicos ejecutados (EDA, particiones, objetivo,
línea base, latencia, secuencias, corrección del \emph{score} adverso y especialista).
\item [\texttt{scripts/}] Puntos de entrada reproducibles de cada experimento.
\item [\texttt{src/}] Paquete con la lógica compartida (datos, características, modelos,
evaluación), con pruebas automatizadas.
\item [\texttt{results/}] Tablas de resultados y ledger anonimizado de 754 unidades de acción.
\item [\texttt{sql/}] Esquema de la consulta base de entrenamiento.
\item [\texttt{reports/}] Esta memoria (fuente LaTeX).
\end{description}

\subsection{Cómo reproducir}

El entorno se prepara con un entorno virtual de Python y la instalación de las dependencias
declaradas. Las pruebas automatizadas verifican la lógica de política y sus guardas. El ledger
\texttt{results/final\_candidate\_actions\_anonymized.csv} permite recalcular las tablas
agregadas, los desgloses por régimen y la incertidumbre agrupada sin acceder al corpus
completo. El script que lo genera queda versionado para mantener trazabilidad, aunque su
ejecución integral requiere el workspace privado de investigación.

\subsection{Datos no publicados}

El conjunto completo (del orden de 2,6 millones de filas \emph{cross-venue}, más la base de
captura) no se publica por tamaño. En su lugar se incluyen el esquema SQL, un ledger
anonimizado suficiente para reproducir las cifras finales y el código del \emph{pipeline}.
Esto permite una auditoría parcial, no la reproducción bit a bit del entrenamiento. Declarar
esa frontera de forma explícita es parte de la trazabilidad y reduce la deuda técnica propia de
estos sistemas~\cite{sculley2015debt}.
